% Round 17. Author: vector_boundary branch.
% Status: CANDIDATE pending independent review.
% Depends on CD17 curve comparison and KC17 compatible-kernel duality.
\section{A vector boundary criterion on a forest of components}
\label{r17ft:section}

The numerical two-component condition of CD17 has a many-component
extension which retains the entire restriction tuple.  It is not
an edgewise replacement for the total-image condition in GG16.
The new input is that the graph supplies all positive directions
of the sum of the divisor Hodge forms when its dual graph is a
forest.

\subsection{The simultaneous boundary spaces}

Let $\mathcal X/R$ be integral, regular, projective and reduced
strictly semistable of relative dimension $n\geq3$.  Suppose its
generic fiber is smooth geometrically integral and that every
special-fiber component $D_v$ and every connected double stratum
$E_e$ is smooth geometrically integral over the native finite
residue field $k$.  Let $\mathcal G$ be the finite multigraph
having those components as vertices and those double strata as
edges, and suppose $\mathcal G$ is a forest.  In particular there
are no triple strata: a triple point would give a triangle in
this multigraph.  Orient every edge, and write $e^-,e^+$ for
its tail and head and $j_{ev}:E_e\hookrightarrow D_v$ for an
incident inclusion.

Fix a relative ample divisor $H$ on the model.  Its restrictions
are denoted $H_v$ and $h_e$.  For $a\geq0$ let
\[
 V_a=\bigoplus_vN^a(D_v)_{\mathbf Q},\qquad
 W_a=\bigoplus_eN^a(E_e)_{\mathbf Q},\qquad
 (\rho_a u)_e=j_{e,e^+}^*u_{e^+}-j_{e,e^-}^*u_{e^-}.
 \tag{FT17.1}\label{r17ft:rho}
\]
Define the simultaneous spaces
\[
                 R=\operatorname{im}\rho_1\subset W_1,
       \qquad C=\operatorname{im}\rho_{n-2}\subset W_{n-2}.
 \tag{FT17.2}\label{r17ft:spaces}
\]
All groups in \eqref{r17ft:rho}--\eqref{r17ft:spaces} are
numerical.  On $W_1$ use the rational symmetric form
\[
       B(u,v)=\sum_e\deg_k(u_ev_eh_e^{n-3}).
 \tag{FT17.3}\label{r17ft:form}
\]
The notation $h^{n-3}R$ means componentwise multiplication by
the specified $h_e^{n-3}$.

\begin{lemma}[The forest fills the positive directions]
\label{r17ft:positive}
For each edge $e$, the vector equal to $h_e$ on that edge and
zero on every other edge belongs to $R$.  Consequently $B|_R$
is nondegenerate.
\end{lemma}

\begin{proof}
Removing $e$ disconnects its tree into two parts.  Put $H_v$
on all vertices of one part and zero on all other vertices.
Every edge besides $e$ has equal scalar coefficients at its
two ends, so its restriction difference is zero.  The difference
on $e$ is $h_e$ or $-h_e$.  Changing the sign gives the stated
vector.  The construction is performed in numerical groups,
with a common relative ample class, so it does not posit
independent algebraic polarizations on intersecting components.

For every edge, the divisor Hodge index theorem gives signature
$(1,\rho(E_e)-1)$ to $\deg(uvh_e^{n-3})$ and makes $h_e$ a
positive vector.  The preceding edge vectors therefore span
a positive definite subspace $P\subset R$ of dimension equal
to the number of edges.  Its orthogonal complement in $W_1$
is negative definite.  Orthogonal projection onto $P$ uses a
diagonal matrix of the nonzero rational numbers
$\deg h_e^{n-1}$; hence it preserves $R$.  Thus
\[
                    R=P\oplus(R\cap P^\perp),
\]
with positive and negative definite summands.  In particular
the restriction of $B$ to $R$ has no radical.
\end{proof}

\begin{proposition}[Exact vector boundary quotient on a forest]
\label{r17ft:kernel}
Let $A_a:V_a\to V_{a+1}$ be the native numerical incidence
map given by restriction of an actual vertical Chow class.
Then
\[
 C=h^{n-3}R\oplus(C\cap R^\perp),\qquad
 \frac{\ker A_{n-2}}{\ker\rho_{n-2}}
       \simeq C\cap R^\perp
       \simeq\frac{C}{h^{n-3}R}.
 \tag{FT17.4}\label{r17ft:defect}
\]
The orthogonal complement uses the native degree pairing
$W_{n-2}\times W_1\to\mathbf Q$.
In particular
\[
          \ker A_{n-2}=\ker\rho_{n-2}
          \quad\Longleftrightarrow\quad C=h^{n-3}R.
 \tag{FT17.5}\label{r17ft:equivalence}
\]
\end{proposition}

\begin{proof}
The reduced total-fiber relation gives the diagonal
self-intersection term on $D_v$ as minus the sum of its
incident double divisors.  The off-diagonal term is the
Gysin push of the restriction from the adjacent component.
With the chosen orientation, this implies
\[
 \langle A_{n-2}u,w\rangle
       =-\sum_e\deg_k\bigl((\rho_{n-2}u)_e(\rho_1w)_e\bigr)
               \qquad(u\in V_{n-2},\ w\in V_1).
 \tag{FT17.6}\label{r17ft:incidence}
\]
Here the component pairing is the unweighted native pairing
of complementary numerical classes.  Its nondegeneracy
shows that $A_{n-2}u=0$ if and only if
$\rho_{n-2}u\in C\cap R^\perp$.  Since $\rho_{n-2}$ has
image $C$, it induces the first quotient identification.

For every $w\in V_1$ one has
\[
          \rho_{n-2}((H_v^{n-3}w_v)_v)
                    =h^{n-3}\rho_1(w).
\]
Thus $h^{n-3}R\subset C$.  Given $c\in C$, the functional
$v\mapsto\sum_e\deg(c_ev_e)$ on $R$ has a unique representing
element $u\in R$ under the nondegenerate rational form of
Lemma~\ref{r17ft:positive}.  Then $c-h^{n-3}u$ belongs to
$C\cap R^\perp$.  The same nondegeneracy makes the two
summands disjoint.  This proves the decomposition and the
remaining assertions.  Every construction is over $\mathbf Q$.
\end{proof}

\subsection{Application to actual strong Green forms}

For every double stratum let $\mathcal T_{E_e}$ denote the
curve-class defect of CD17: the kernel of the numerical
projection from the native algebraic span in
$H^{2n-4}(E_{e,\bar k},\mathbf Q_\ell(n-2))$.

\begin{theorem}[Forest codimension-two correction]
\label{r17ft:theorem}
Under the model hypotheses above, suppose for one $\ell\ne p$
that
\[
           \mathcal T_{E_e}=0\quad\hbox{for every }e,
                    \qquad C=h^{n-3}R.
 \tag{FT17.7}\label{r17ft:conditions}
\]
Then every geometrically homologically trivial rational
codimension-two class on the generic fiber has a numerically
admissible rational Chow extension to $\mathcal X$.
Consequently each specified such cycle has an ordinary actual
strong Green form on its exact support complement, with every
compact actual test and the original valued-field normalization.
\end{theorem}

\begin{proof}
Proposition~\ref{r17ft:kernel} gives
$\ker A_{n-2}=\ker\rho_{n-2}$.
For a tuple in this numerical kernel choose arbitrary actual
rational Chow representatives on the components.  Their
restriction difference on each $E_e$ is a numerically zero
curve cycle.  Its geometric class is zero by the assumed
curve comparison.  The tuple is therefore cohomologically
compatible on every double stratum.

This is exactly the reduced compatible-kernel hypothesis of
KC17, Proposition \texttt{r17kc:criterion}, for $r=2$.
That proposition lifts the compatible cohomology tuple to
the special fiber and then to the strict-henselian model;
cup product with the global pure operational class of a
generic homologically trivial Chow extension has zero native
weighted top trace.  The restriction tuple thus annihilates
$\ker A_{n-2}$.  Perfect rational numerical duality and
$A_1^{\mathsf T}=A_{n-2}$ produce a rational vertical divisor
correction whose component restrictions cancel it numerically.
This proves the stated admissibility assertion.

Apply the frozen pure-$K_0$ model witness and native actual
dominant detector to the corrected extension.  Its actual
obstruction vanishes.  Reconstruction from the curved pair
for the original specified cycle, rather than a new moved
cycle, gives the claimed exact-support strong G identity.
This last step retains all actual compact tests and the full
native degrees, precisely as in BT17 and CD17.
\end{proof}

\begin{corollary}[One-sided divisor lifting along a rooted tree]
\label{r17ft:rooted}
Orient each tree away from a chosen root.  Suppose every
nonroot component $D_v$ restricts surjectively onto the
numerical divisor group of its parent double stratum:
\[
                  N^1(D_v)_{\mathbf Q}
                         \twoheadrightarrow N^1(E_{v,\mathrm{parent}})_{\mathbf Q}.
 \tag{FT17.8}\label{r17ft:one-sided}
\]
Then $R=W_1$ and $C=W_{n-2}$, so the numerical condition in
\eqref{r17ft:conditions} is automatic.  If each double
stratum has curve comparison, all conclusions of
Theorem~\ref{r17ft:theorem} hold.

In particular it is sufficient that each double stratum be
a projective bundle over a smooth projective variety of
dimension at most two, or a specified smooth modification
or generically finite dominated variety covered by CD17.
No divisor Tate hypothesis on those base surfaces is needed.
\end{corollary}

\begin{proof}
Give an arbitrary desired tuple $t_e\in N^1(E_e)$.
Set the class at every root to zero.  Proceeding away from
the root, after the parent class $u_p$ has been fixed, use
\eqref{r17ft:one-sided} to choose $u_v$ satisfying
\[
           j_{e,v}^*u_v=t_e+j_{e,p}^*u_p.
\]
The restriction difference on the parent edge is then exactly
$t_e$.  Its values on the other incident edges impose no
previous condition: the corresponding child classes are
chosen at subsequent steps.  Finiteness and absence of cycles
therefore produce a tuple $u$ with $\rho_1u=t$.
This proves $R=W_1$ and retains all simultaneous edge data.

Componentwise numerical divisor Lefschetz identifies
$h^{n-3}W_1$ with $W_{n-2}$.  Since
$h^{n-3}R\subset C\subset W_{n-2}$, this also gives
$C=W_{n-2}$.  The theorem and the proved curve-comparison
families of CD17 now apply.
\end{proof}

\paragraph{Beyond a forest.}
The kernel identity \eqref{r17ft:incidence} holds for an
arbitrary reduced strictly semistable model.  Thus full
simultaneous surjectivity of $\rho_1$ together with curve
comparison on all doubles already suffices, without a
forest assumption, by the same compatible-kernel proof.
However the forest's nondegeneracy argument need not hold
for a graph with cycles: the scalar ample differences then
give only the cut space of the graph.  We do not replace
surjectivity by independent edgewise surjectivity or infer
that a graph cycle contributes no obstruction.

\paragraph{Scope and historical distinction.}
The essential new conditions are curve-class comparison
and a simultaneous numerical divisor test space.  The
proof does not choose algebraic preimages of the supported
cohomology residuals, does not assume divisor Tate, and
does not reduce a vector of boundary images to scalar
edge coefficients.  The signed-preimage and scalar graph
criteria of R7 and R16 remain valid in their stated scopes.
No claim is made that an arbitrary variety admits a model
satisfying \eqref{r17ft:conditions} or \eqref{r17ft:one-sided}.
For fourfolds the known endpoint and one-cycle cases turn
this codimension-two theorem into full local G.  A complete
number-field height still requires the local hypothesis
at every finite place.
